\chapter{Architecture logicielle et reproductibilité}
\label{chap:architecture}

Ce chapitre décrit l'organisation du prototype afin de distinguer les choix scientifiques des détails nécessaires à leur reproduction. Le dépôt associe le code source, les configurations et la documentation. Les résultats lourds, tels que les modèles entraînés, animations et tableaux bruts, sont produits séparément à partir de ces éléments versionnés.

\section{Organisation des modules}
Le simulateur est séparé des politiques de contrôle. Il prend en charge l'évolution des positions, le calcul du voisinage, les limites de l'arène et les pannes. Cette séparation permet d'exécuter une même condition initiale avec une règle de regroupement ou avec le GNN, sans modifier les mécanismes physiques simplifiés. Les générateurs de formes transforment ensuite les images géométriques et symboliques en nuages de points normalisés. Enfin, les modules d'entraînement et d'évaluation construisent les lots d'épisodes, calculent les pertes et sauvegardent les mesures.

Cette architecture évite qu'une fonction destinée à la visualisation influence involontairement la dynamique de simulation. Elle rend aussi les tests unitaires plus ciblés : le voisinage, le nombre variable de robots, l'invariance à la permutation et le traitement d'un robot isolé peuvent être testés indépendamment de la qualité finale d'une forme.

\section{Cycle d'exécution d'un épisode}
Un épisode commence par le tirage de positions initiales et par le choix d'une forme cible. À chaque pas, le simulateur construit le graphe de communication à partir des positions courantes. Chaque robot actif encode alors son observation locale, agrège les messages de ses voisins et produit une vitesse bornée. Le simulateur applique simultanément les commandes, puis la boucle recommence. Après le dernier pas, la configuration finale est comparée au nuage de points de la cible.

Pendant l'entraînement, cette boucle est différentiable et la perte remonte jusqu'aux paramètres du GNN. Pendant l'évaluation, les paramètres sont gelés et le même mécanisme ne produit que les métriques. Cette séparation réduit le risque d'évaluer un modèle dans un mode différent de celui utilisé pour son entraînement.

\section{Configurations et traçabilité}
Les paramètres numériques ne sont pas dispersés dans le code : ils sont regroupés dans des fichiers de configuration. Les principaux fichiers fixent notamment la graine aléatoire, la taille de lot, le nombre d'itérations, le rayon de communication, la distance de sécurité, les poids de la perte et les étapes du curriculum. Chaque exécution sauvegarde la configuration utilisée avec les courbes et mesures produites.

La graine ne rend pas toute expérience scientifique définitive ; elle permet en revanche de reproduire exactement un épisode et d'identifier les divergences lorsqu'une implémentation évolue. L'étude finale doit donc conserver à la fois les graines individuelles et les agrégats sur un nombre plus élevé de répétitions.

\section{Visualisation et inspection}
Les figures de formation ne sont pas seulement des éléments de présentation. Elles servent à détecter des comportements qui peuvent échapper à une moyenne numérique : amas local, points de cible non couverts, trajectoires instables ou collisions récurrentes. Les images de fin d'épisode et les animations sont donc générées à partir du même simulateur que celui utilisé dans les métriques.

L'approche conserve toutefois une frontière nette entre illustration et preuve. Une animation convaincante pour une graine ne remplace pas une évaluation multi-graines ; inversement, une moyenne faible doit être confrontée à la distribution spatiale réellement obtenue. Le chapitre suivant applique ce principe aux résultats disponibles.
